\subsection{Duration Matching Ablation}

Feature extraction for the classical baselines uses a fixed audio duration per utterance. Table~\ref{tab:duration_ablation} compares three duration configurations, evaluated using the selected WavLM-large, all-layer-averaged representation on the same 20,000-utterance development subset: a fixed 1.0-second duration, a duration sampled uniformly at random in the 1.0--3.0-second range per utterance, and a fixed 3.0-second duration.

\renewcommand{\arraystretch}{1.4}
\setlength{\tabcolsep}{8pt}
\begin{table}[H]
\centering
\caption{Duration Configuration Ablation (Development Set EER, \%)}
\label{tab:duration_ablation}
\begin{tabularx}{\textwidth}{>{\raggedright\arraybackslash}p{3.5cm} *{4}{>{\centering\arraybackslash}X}}
\hline
\textbf{Duration Configuration} & \textbf{XGBoost} & \textbf{Random Forest} & \textbf{SVM} & \textbf{GBM} \\
\hline
Fixed 1.0s & 16.01 & 16.96 & 10.14 & 13.54 \\
Random 1.0--3.0s & 10.19 & 10.98 & 6.23 & 7.74 \\
Fixed 3.0s (selected) & 6.66 & 7.13 & 2.86 & 5.38 \\
\hline
\end{tabularx}
\end{table}

The fixed 3.0-second configuration outperformed both alternatives across all four models, improving EER by 4--6 points over the random 1.0--3.0s range and by 8--10 points over the fixed 1.0-second configuration. A consistent, longer analysis window benefits classical feature extraction more than duration diversity matching the anticipated 1--3 second inference range. Short, single-second windows appear to cut out spectro-temporal structure that matters for spoof detection. The fixed 3.0-second configuration was adopted for all subsequent classical-baseline results in this section.